\lyseis{Α1.0}{Εισαγωγή: οι φυσικοί αριθμοί και οι πράξεις τους}
\sol{1} Ο $0$ δεν έχει προηγούμενο και έχει επόμενο τον $1$. Ο $99$: $98$ και $100$. Ο $1.000$: $999$ και $1.001$. Ο $25.099$: $25.098$ και $25.100$.
\sol{2} Από αριστερά: $4$ εκατομμύρια, $6$ εκατοντάδες χιλιάδων, $3$ δεκάδες χιλιάδων, $5$ χιλιάδες, $1$ εκατοντάδα, $8$ δεκάδες, $2$ μονάδες. Το $6$ είναι στις εκατοντάδες χιλιάδων: αξία $600.000$. Το $3$ είναι στις δεκάδες χιλιάδων: αξία $30.000$.
\sol{3} α) $<$ β) $>$ (ίδιο πλήθος ψηφίων, τα τρία πρώτα ίσα και στη συνέχεια $7>0$) γ) $>$ (ο πρώτος έχει περισσότερα ψηφία) δ) $=$.
\sol{4} α) αντιμεταθετική της πρόσθεσης β) προσεταιριστική της πρόσθεσης γ) το $1$ είναι ουδέτερο στοιχείο του πολλαπλασιασμού δ) το $0$ είναι απορροφητικό στοιχείο ε) επιμεριστική.
\sol{5} α) $7\cdot(100+2)=700+14=714$. β) $8\cdot(100-1)=800-8=792$. γ) $25\cdot(10+2)=250+50=300$.
\sol{6} Πρώτος τρόπος: $6\cdot 25+4\cdot 25=150+100=250$ κιλά. Δεύτερος τρόπος: $(6+4)\cdot 25=10\cdot 25=250$ κιλά. Είναι η επιμεριστική ιδιότητα.
\sol{7} α) $40-30+8=18$. β) $35\cdot 6+8=210+8=218$. γ) $2\cdot 7+10=14+10=24$. δ) $3+8\cdot 5=3+40=43$.
\sol{8} Ο μικρότερος τετραψήφιος είναι ο $1.000$ και ο μεγαλύτερος τριψήφιος ο $999$. Οι διψήφιοι είναι από $10$ ως $99$, δηλαδή $99-10+1=90$. Οι άρτιοι είναι οι $10,12,\dots,98$, δηλαδή $45$ (οι μισοί).
\sol{9} Έχουμε $8-5=3$, ενώ το $5-8$ δεν υπάρχει στους φυσικούς. Άρα η αντιμεταθετική δεν ισχύει. Επίσης $(8-3)-2=5-2=3$, ενώ $8-(3-2)=8-1=7$: διαφέρουν, άρα δεν ισχύει ούτε η προσεταιριστική.
\sol{10} Μεγαλύτερος: $8.520$. Μικρότερος: $2.058$ (το $0$ δεν μπορεί να μπει πρώτο). Διαφορά: $8.520-2.058=6.462$.
